% !PRINTED 542
\weberpage{542}

\begin{center}
\S.\,164.\\
\gesperrt{Resolventen von Lagrange.}
\end{center}

Die Methode der Auflösung cyklischer Gleichungen, die wir jetzt kennen lernen wollen, ist gleichmässig auf Primzahlgrade und zusammengesetzte Grade anwendbar. Man bedient sich dazu gewisser Ausdrücke, die unter dem Namen der \gesperrt{Resolventen von Lagrange} bekannt sind\footnote{\gesperrt{Lagrange}, Réflexions etc., s. S. 508. Früher haben wir unter Resolventen auflösende Gleichungen verstanden, hier sind es auflösende Functionen.}, die bei allen Untersuchungen über die algebraische Auflösung von Gleichungen von grossem Nutzen sind.

Es sei $F(x) = 0$ eine Gleichung mit den Wurzeln $\alpha, \alpha_1, \alpha_2 \ldots \alpha_{m-1}$. Wir bezeichnen mit $\varepsilon$ irgend eine $m^{\text{te}}$ Einheitswurzel und führen die Bezeichnung ein
\weq{1}{(\varepsilon, \alpha) = \alpha + \varepsilon\alpha_1 + \varepsilon^2\alpha_2 + \cdots + \varepsilon^{m-1}\alpha_{m-1}.}
Die so definirten Summen sind es, die man die \gesperrt{Lagrange'schen Resolventen} nennt. Wenn diese Functionen für alle $m^{\text{ten}}$ Einheitswurzeln $\varepsilon$ bekannt sind, so ist auch die Gleichung selbst gelöst, denn es ist nach \S.\,133 (6)
\weq{2}{\overset{\varepsilon}{\Sigma}\,\varepsilon^k = m \ \text{oder} = 0,}
je nachdem $k$ durch $m$ theilbar ist oder nicht, und daraus folgt
\weq{3}{m\alpha = \overset{\varepsilon}{\Sigma}(\varepsilon, \alpha),}
worin sich die Summen über alle $m^{\text{ten}}$ Einheitswurzeln $\varepsilon$ erstrecken. Man kann auch die anderen Wurzeln in gleicher Weise ausdrücken:
\weq{4}{m\alpha_k = \overset{\varepsilon}{\Sigma}\varepsilon^{-k}(\varepsilon, \alpha),}
so dass in der That Alles auf die Kenntniss von $\varepsilon$ und der Functionen $(\varepsilon, \alpha)$ zurückgeführt ist.

Die Summen (2) und (3) lassen sich noch in etwas anderer Weise darstellen, da die sämmtlichen $m^{\text{ten}}$ Einheitswurzeln Potenzen einer primitiven unter ihnen sind. Versteht man also unter $\varepsilon$ eine festgehaltene primitive $m^{\text{te}}$ Einheitswurzel und unter $\lambda$ einen Index, der ein volles Restsystem nach dem Modul $m$, etwa die
%CONT>>
